%!TEX root = ../template.tex
%%%%%%%%%%%%%%%%%%%%%%%%%%%%%%%%%%%%%%%%%%%%%%%%%%%%%%%%%%%%%%%%%%%%
%% chapter3.tex
%% NOVA thesis document file
%%
%% Chapter 3: Related Work
%%%%%%%%%%%%%%%%%%%%%%%%%%%%%%%%%%%%%%%%%%%%%%%%%%%%%%%%%%%%%%%%%%%%

\typeout{NT FILE chapter3.tex}%

\chapter{Related Work}
\label{cha:related_work}

\section{Conceptual modelling education landscape}
\label{sec:rw_landscape}

Research on conceptual modelling education spans multiple modelling families and research communities, including data modelling, process modelling, requirements-related modelling, and model-driven development. Maslov et al.\ map the conceptual modelling education literature and report that educational work is distributed across multiple communities and modelling subfields that have developed partially independent terminology, learning-theory grounding, and evaluation approaches \cite{maslov2025conceptual}. As a result, findings are often difficult to compare across modelling families, and evidence about effective teaching strategies and supports accumulates unevenly across subdomains.

This thesis builds on the landscape evidence by drawing together frameworks and empirical results from different modelling-education strands (conceptual/data, process, goal-oriented, and tool-supported modelling) and by focusing on the challenges and support needs reported by both students and instructors.

\section{Modelling frameworks}
\label{sec:rw_frameworks}

A common approach to strengthening modelling education research is to define explicit frameworks for learning outcomes, competencies, and assessment design. Several works adapt Bloom's revised taxonomy \cite{krathwohl2002revision} to modelling activities in order to clarify what learners should know and be able to do beyond learning a notation. The following subsections summarise three modelling-education frameworks that are especially relevant to this thesis.

\subsection{CaMeLOT: a framework for conceptual data modelling learning outcomes}
\label{subsec:rw_camelot}

CaMeLOT proposes an educational framework for conceptual data modelling that combines Bloom's cognitive levels with a knowledge dimension (factual, conceptual, procedural, metacognitive) and modelling content areas \cite{bogdanova2019camelot}. The framework is instantiated with detailed example learning outcomes and can be used to design course objectives, learning activities, and assessments in a systematic manner. For this thesis, CaMeLOT is valuable as a template for articulating modelling competence across cognitive and knowledge dimensions, and for motivating the inclusion of procedural and reflective elements rather than focusing exclusively on diagram production.

\subsection{Domain Modelling in Bloom: empirical analysis of assessment practices}
\label{subsec:rw_domain_bloom}

Bogdanova and Snoeck empirically analyse domain-modelling assessment tasks through the lens of Bloom's revised taxonomy \cite{bogdanova2017domain}. They report that many tasks concentrate on understanding, analysing, and creating conceptual knowledge, while evaluation, procedural knowledge, and metacognitive learning outcomes receive limited explicit attention. This work provides evidence that modelling courses may implicitly expect higher-level modelling competence without systematically scaffolding or assessing foundational skills. For this thesis, these findings support examining modelling challenges in relation to cognitive demands and considering whether educational support should target not only syntax and notation use but also strategies for evaluation and refinement.

\subsection{Teaching conceptual modelling and metamodeling}
\label{subsec:rw_bork_framework}

Bork extends Bloom-based framing to both conceptual modelling and metamodeling by relating cognitive and knowledge dimensions to components of modelling methods (language, semantics, notation, procedure, mechanisms and algorithms) \cite{bork2019framework}. By analysing a teaching case using the proposed framework, the work suggests that even practice-oriented materials may emphasise model creation while providing limited explicit support for evaluation tasks and metacognitive reflection. This perspective is useful for the present thesis because it frames modelling competence as method-level understanding and reflective judgement, aligning with the thesis focus on the transition from notation use to meaningful modelling practice.

\section{Empirical studies on modelling learning challenges}
\label{sec:rw_empirical_challenges}

Frameworks clarify educational targets, but empirical studies are needed to understand where learners struggle in practice and which supports are perceived as valuable. Two complementary strands are particularly relevant: studies of goal-oriented modelling learning challenges and research agendas for process modelling education.

\subsection{iStar learning challenges and support requirements}
\label{subsec:rw_istar_learning}

A key empirical reference for this thesis is an iStar learning study, which investigates challenges beginners face when learning and practising iStar and derives requirements for facilitating learning \cite{istar-learning}. Using Bloom-based task framing and grounded theory analysis, the study reports difficulties at multiple cognitive levels, including understanding the semantics of modelling constructs, applying the notation to requirements scenarios, and analysing models for quality and completeness. The authors derive a catalogue of strategies and tool-supported requirements (e.g., conceptual clarification, syntax checking, process guidance, templates, advanced visualisation, and partial automation from textual requirements) and implement a subset of these requirements in a prototype tool \cite{istar-learning}. This thesis builds on these insights by extending the focus from a single modelling method to broader modelling education contexts and by eliciting perspectives from both students and instructors.

\subsection{Process modelling education and the need for empirical validation}
\label{subsec:rw_bpmn_education}

In the domain of process modelling, Maslov argues for empirically validated process-modelling education using BPMN and outlines a research agenda that integrates conceptual modelling education, process model quality research, and learning theory \cite{bpmn-education}. The agenda emphasises systematic synthesis of existing work, empirical investigation of novice modelling behaviour and typical errors, and the design and experimental validation of BPMN-specific teaching strategies and feedback mechanisms. This agenda is relevant to the present thesis as a parallel strand: it highlights that modelling families pose distinct cognitive and practical demands while pointing to recurring educational needs such as scaffolding, feedback, and support for managing modelling complexity.

\section{Tools and model-driven engineering in education}
\label{sec:rw_tools_mde}

Modelling education is strongly shaped by tool choice, especially in model-driven and model-based engineering settings. Tools can support learners through feedback, automation, and manageability, but they can also introduce friction that influences how students perceive modelling tasks and their own competence.

\subsection{Perception of tools and UML in MDE education}
\label{subsec:rw_liebel_tools_uml}

Liebel et al.\ show that students’ tool experiences depend less on the tool “in isolation” than on how it is embedded in the course---including the support provided, the surrounding workflow, and the expectations placed on models (for example, exploratory communication artefacts versus precise inputs to downstream automation) \cite{liebel2017model}. The study also suggests that perceptions of UML are more robustly positive when models are used for requirements and analysis rather than primarily as code generation inputs \cite{liebel2017model}. These results motivate treating tooling and course design as socio-technical factors that shape modelling experiences, rather than as neutral delivery mechanisms.

\subsection{Human factors in model-driven engineering}
\label{subsec:rw_hf_mde}

Work on human factors in model-driven engineering argues for research that accounts for usability, cognitive load, and the alignment between tools and user goals \cite{human-factors-mde}. From an educational perspective, this supports the view that modelling tools should be evaluated not only by feature completeness, but also by how they support learning processes such as exploration, error detection, reflection, and progressive mastery.

\section{Human factors and learning context in modelling education}
\label{sec:rw_human_context}

Beyond frameworks, methods, and tools, modelling education is influenced by motivational and contextual factors such as domain choice, inclusiveness, and the authenticity of tasks.

\subsection{Domain diversity, motivation, and inclusion}
\label{subsec:rw_domain_diversity}

Research on domain diversity and motivation in modelling education reports that educators' assumptions about ``motivating'' domains do not always match students' preferences, and that students value domains that connect to their interests and everyday or socially relevant contexts \cite{grassl2026dei}. The work also highlights that assignment design can unintentionally exclude learners when domains or examples presume cultural, economic, or experiential background that not all students share \cite{grassl2026dei}. These findings motivate treating domain choice and inclusiveness as relevant contextual variables when interpreting modelling learning experiences.

\subsection{Experiential learning and authenticity}
\label{subsec:rw_experientialism}

Holmes et al.\ propose dimensions of experientialism in software engineering education and show that students value learning experiences that involve real projects, real tasks, and real mentors \cite{holmes2018dimensions}. Although this work is not modelling-specific, it contributes a useful lens for interpreting how authenticity and mentorship can influence learners' engagement and competence development, and it supports the relevance of studying modelling experiences in realistic educational contexts.

\section{Competence-oriented agendas and methodological precedents}
\label{sec:rw_competence_methods}

Competence-oriented perspectives provide a complementary lens for structuring educational goals and for linking teaching approaches to measurable outcomes.

\subsection{Competency identification and development in SE education}
\label{subsec:rw_evelin}

Sedelmaier and Landes propose the EVELIN research agenda for identifying and developing competencies in software engineering and argue that curricula should be grounded in explicit competence models rather than topic lists \cite{sedelmaier2013research}. They introduce a six-level classification system for technical competencies and advocate an iterative approach in which competence targets are identified, competence levels are measured, and teaching strategies and structural constraints are adjusted over time \cite{sedelmaier2013research}. While not specific to modelling, this agenda supports treating modelling competence as multi-dimensional (conceptual, procedural, analytical, socio-technical) and complements Bloom-based modelling frameworks by emphasising explicit competence models and iterative alignment between competence targets, teaching methods, and structural constraints.

\subsection{Grounded Theory in software engineering research}
\label{subsec:rw_gt}

Methodologically, grounded theory has been widely discussed in software engineering research as an approach for developing explanatory theory from qualitative data. Stol et al.\ provide a critical review and practical guidelines for conducting grounded theory in software engineering \cite{stol2016grounded}, while Hoda et al.\ propose a socio-technical grounded theory perspective that explicitly accounts for the interplay between social and technical factors \cite{hoda2018socio}. Adolph et al.\ offer a practical account of using grounded theory to study the experience of software development \cite{adolph2011using}. These works provide methodological grounding for the approach used in this thesis and support the use of interviews and iterative coding to derive an explanatory account from participant experiences.

\section{Relating the literature to the present study}
\label{sec:rw_dialogue}

This section places the literature above, and a small set of additional papers, against the 14-category account in Chapter~\ref{cha:findings}. The additional papers were read because they speak to claims that the core set covers only thinly---delayed critique of student diagrams, group versus individual UML work, consistency across modelling views, and the reliability of LLM-generated UML. They strengthen selected claims; they do not replace the core literature.

\subsection{How the core literature speaks to the findings}
\label{subsec:rw_core_against_findings}

Maslov et al.'s picture of conceptual modelling education as fragmented across families \cite{maslov2025conceptual} is the landscape condition under which the present corpus was collected: interviews span UML, BPMN, ER, iStar, and courses that teach several of those notations together, rather than a single-language laboratory. The same fragmentation is the practical cost later labelled CAT14.

The Bloom-informed frameworks \cite{bogdanova2019camelot,bogdanova2017domain,bork2019framework} name an evaluation and metacognitive gap in intended learning outcomes. That vocabulary is useful for CAT12 and PH2: students invent checklists and wait for the teacher because a compile-like signal is missing. The frameworks describe intended learning outcomes. They do not, by themselves, describe the lived substitutes observed in the interviews, and this study did not sort transcripts into Bloom bins (Chapter~\ref{cha:methodology}).

The closest qualitative peer remains the iStar learning study \cite{istar-learning}. It supports the first-modelling-move problem (PH1) and parts of CAT1, CAT4, and CAT7: moving from a requirements scenario to a first model is hard, and learners ask for explanation and checking rather than only a symbol list. It is also a methods peer (grounded analysis in modelling education). It is not a template that this thesis replicates. Li et al.\ studied one goal-modelling language after a tutorial and a modelling session; this corpus is multi-notation, interview-only, and includes lecturers.

Chakraborty and Liebel report that students often do not understand what a model ``says,'' and that perceptions of modelling polarise between valuable instrument and extra documentation \cite{chakraborty2022perceptions}. That finding supports CAT2 as a comprehension and notation problem rather than only a missing demonstration. It does not split the start-gate (CAT1) from overload (CAT2); the present analysis treats those as distinct mechanisms.

Maslov's BPMN education agenda \cite{bpmn-education} remains a parallel strand rather than a source of BPMN learning outcomes. It flags scaffolding, feedback, and typical novice errors as open needs. This thesis does not close that agenda; it shows how getting started, getting comments, using tools, and working in groups hang together across families.

Liebel et al.\ and the human-factors MDE literature \cite{liebel2017model,human-factors-mde} support treating tools as socio-technical rather than as a brand choice. CAT7 versus CAT8 is that tension in interview form: semantic checks are wanted; an antique compile UI or a mid-course licence switch is still dropped. Buchmann et al.\ insist that models are knowledge structures, not drawings \cite{buchmann2019conceptual}, which helps explain why a drawing tool that cannot catch an invalid model is a poor substitute for a checker---and why students still choose draw.io when the checker costs more than it saves.

Grassl et al.\ treat motivating domains and inclusion as design variables \cite{grassl2026dei}; Holmes et al.\ describe authenticity as real projects, real tasks, and real mentors \cite{holmes2018dimensions}. CAT10 relates to that authenticity literature and qualifies it: in these interviews a concrete reader (supervisor, stakeholders, defence) pulled care up more reliably than domain excitement alone. Modelling-specific evidence for ``who reads the diagram'' remains thinner than evidence for ``which domain the diagram is about.''

Course load and theory--practice gaps (CAT11) have little dedicated modelling-education literature. EVELIN treats structural conditions as something competence-oriented curricula must measure \cite{sedelmaier2013research}; that is an adjacent agenda, not a study of ECTS versus a paid job versus packed concurrent courses.

\subsection{Additional work that strengthens selected claims}
\label{subsec:rw_additional}

A further set of papers is useful where the core frame is thin.

Industrial modelling tools give little quality feedback, so students wait on the instructor. Hasker and Rowe built UMLint because a student may work for days, submit, and wait further days for a graded response, while producing diagrams that are syntactically acceptable and semantically wrong \cite{hasker2011umlint}. Foss, Urazova, and Lawrence similarly treat delayed marking of design diagrams as a teaching bottleneck and offer immediate, automated feedback on UML database-design questions \cite{foss2022autoer}. Those papers support CAT5 and PH2 as a recognised educational problem: the correctness signal is late and person-bound unless a tool intercepts it. They also show that automated checkers are an available design response, not only ``more lecturer hours.'' They do not address CAT6. Access---who can show a half-finished diagram, in public or at the desk, on graded or ungraded work---is largely absent from that tool literature.

Silva, Gadelha, Steinmacher, and Conte compared individual and group UML exercises and found that groups were not uniformly better: some artefacts were more correct when produced individually and more complete when produced in a group \cite{silva2018group}. That supports CAT9's two-sided claim---groups can kill or save a diagram---and argues against treating group work as automatically the better pedagogy. Their outcome measures are correctness and completeness of classroom exercises, not code-first pressure, exclusive notation ownership, or a defence audience. Those conditions are the present study's addition.

Sikkel and Daneva argued that textbooks rarely teach consistency across UML views and developed pairwise-comparison exercises so that students notice missing elements between diagrams \cite{sikkel2010consistency}. Verbruggen and Snoeck observed students modelling the same case in a UML class diagram and a BPMN process model and argued that knowing which requirements belong in which perspective is part of model quality, not an extra trick \cite{verbruggen2022multiperspective}. Those sources support CAT14 as a recognised teaching problem: a class diagram, a process model, and a goal model answer different questions about the same system, and keeping them consistent has to be taught. They do not document the incidents in this corpus---an iStar rationale that loses its goals when the work moves to UML (S6), or a BPMN pool copied too directly into classes (S10). They also do not say whether courses teach that difference up front, or whether students only hear it after a diagram has already gone wrong.

C\'amara, Troya, Burgue\~no, and Vallecillo assessed ChatGPT on UML class diagrams with OCL and found limited reliability: syntactic and semantic deficiencies, inconsistency, and poor scaling compared with code generation \cite{camara2023chatgpt}. That result is consistent with CAT13's student stance (explain and compare, do not submit generated homework): if generation is unreliable, using the tool as an explainer is a rational classroom tactic, not only an integrity rule. It also undercuts any claim that current LLMs already close CAT12 by compiling models. The present CAT13 evidence is four student cases, not an evaluation of ChatGPT.

\subsection{Where the literature remains thin}
\label{subsec:rw_remaining_thin}

Taken together, prior work is strongest on landscape fragmentation, Bloom-shaped outcome gaps, iStar-specific challenges, tool friction in MDE courses, delayed diagram feedback, and multi-perspective consistency. It is weaker on feedback \emph{access} (CAT6), groups as kill-or-save under code-first pressure, AI used as an explainer in modelling courses, and course load as a blocker of demonstrations and feedback loops (CAT11). This thesis does not close those gaps by volume of interviews; it offers a bounded, multi-notation, student-and-lecturer account against which that literature can be read.

\section{Summary}
\label{sec:rw_summary_positioning}

In summary, existing work provides: (i) landscape-level evidence that modelling education research is distributed across modelling families and communities \cite{maslov2025conceptual}; (ii) frameworks that structure modelling learning outcomes using Bloom-based dimensions \cite{bogdanova2019camelot,bogdanova2017domain,bork2019framework}; (iii) empirical studies and agendas that identify modelling-specific learning challenges and the need for validated educational support \cite{istar-learning,bpmn-education}; and (iv) tool-, competence-, and human-centred perspectives that emphasise the role of context, usability, motivation, inclusiveness, and authenticity \cite{grassl2026dei,holmes2018dimensions,liebel2017model,sedelmaier2013research}. Section~\ref{sec:rw_dialogue} places that frame against the findings of this thesis and adds a smaller set of studies on delayed diagram feedback, group versus individual UML work, multi-perspective modelling, and LLM-generated UML \cite{camara2023chatgpt,foss2022autoer,hasker2011umlint,silva2018group,sikkel2010consistency,verbruggen2022multiperspective}. Those additional sources strengthen selected claims; they do not replace the core literature summarised above.
